% !TEX program = xelatex
\documentclass[11pt]{article}
\usepackage[UTF8]{ctex}
\usepackage{amsmath, amssymb, amsthm}
\usepackage[margin=2.5cm]{geometry}
\usepackage{enumitem}

\newcounter{problem}
\newenvironment{problem}[1][0]
  {\refstepcounter{problem}\par\medskip\noindent\textbf{第\theproblem 题}（#1 分）\par\nopagebreak\noindent\ignorespaces}
  {\par\medskip}

\title{\textbf{中国科学院大学 数学学科 硕转博资格考试综合笔试}\\[4pt]
       \textbf{《代数学》试卷}}
\author{考试时间：180 分钟\qquad\qquad 满分：150 分}
\date{}

\begin{document}
\maketitle
\begin{center}\rule{\textwidth}{0.4pt}\end{center}

\textbf{注意事项：}1. 答案写在答题纸上并注明题号；2. 写明关键步骤，仅有结论不得分；3. 可使用教材中的标准定理，但须说明其条件成立。

\vspace{1em}

% ===== 一、Galois 理论和域论 =====
\section*{一、Galois 理论和域论（共 30 分）}

\begin{problem}[15]
（在此编写题目）
\end{problem}

\begin{problem}[15]
（在此编写题目）
\end{problem}

% ===== 二、模论（共 30 分） =====
\section*{二、模论（共 30 分）}

\begin{problem}[15]
（在此编写题目）
\end{problem}

\begin{problem}[15]
（在此编写题目）
\end{problem}

% ===== 三、环与代数（共 30 分） =====
\section*{三、环与代数（共 30 分）}

\begin{problem}[15]
（在此编写题目）
\end{problem}

\begin{problem}[15]
（在此编写题目）
\end{problem}

% ===== 四、有限群及其表示论（共 30 分） =====
\section*{四、有限群及其表示论（共 30 分）}

\begin{problem}[15]
（在此编写题目）
\end{problem}

\begin{problem}[15]
（在此编写题目）
\end{problem}

% ===== 五、同调代数（共 30 分） =====
\section*{五、同调代数（共 30 分）}

\begin{problem}[15]
（在此编写题目）
\end{problem}

\begin{problem}[15]
（在此编写题目）
\end{problem}

\end{document}